% (c) Nikita Lisitsa, lisyarus@gmail.com, 2026

\documentclass[10pt]{beamer}

\usepackage[T2A]{fontenc}
\usepackage[russian]{babel}
\usepackage{minted}

\usepackage{graphicx}
\graphicspath{ {./images/} }

\usepackage{adjustbox}

\usepackage{tikz}

\usepackage{color}
\usepackage{soul}

\usepackage{hyperref}

\usetheme{metropolis}
\setminted{fontsize=\footnotesize}

\definecolor{blue}{rgb}{0,0,1}
\definecolor{red}{rgb}{1,0,0}
\definecolor{codebg}{RGB}{29,35,49}

\makeatletter
\newcommand{\slideimage}[1]{
  \begin{figure}
    \begin{adjustbox}{width=\textwidth, totalheight=\textheight-2\baselineskip-2\baselineskip,keepaspectratio}
      \includegraphics{#1}
    \end{adjustbox}
  \end{figure}
}
\makeatother

\title{Компьютерная графика}
\subtitle{Домашнее задание 2: Визуализатор сцены с освещением и тенями}
\date{2026}

\setbeamertemplate{footline}[frame number]

\begin{document}

\frame{\titlepage}

\begin{frame}[fragile]
\frametitle{Задание}
\begin{itemize}
\item Сделать визуализатор сцен в формате OBJ с текстурами и освещением по Фонгу
\pause
\item Путь до сцены задаётся, например, параметром командной строки
\pause
\item Два источника света: направленный (`солнце') и точечный, оба как-то двигаются со временем
\pause
\item От обоих источников света есть тени, построенные алгоритмом shadow mapping (+ PCF или VSM, с размытием)
\pause
\item Камера должна управляться пользователем (например, так же, как в первом домашнем задании)
\end{itemize}
\end{frame}

\begin{frame}[fragile]
\frametitle{Сцена}
Хорошая сцена:
\begin{itemize}
\item Порядка 100k-1kk треугольников в сумме
\item Порядка 100-1000 различных объектов
\item У большинства объектов есть текстура альбедо
\end{itemize}
Рекомендую тестировать на \textit{Crytek Sponza}, её можно найти здесь: \href{https://casual-effects.com/data/index.html}{\nolinkurl{casual-effects.com/data/index.html}}
\end{frame}

\begin{frame}[fragile]
\frametitle{Crytek Sponza}
\slideimage{crytek-sponza.jpg}
\end{frame}

\begin{frame}[fragile]
\frametitle{Crytek Sponza}
\slideimage{sponza_real.jpeg}
\end{frame}

\begin{frame}[fragile]
\frametitle{Crytek Sponza}
\begin{itemize}
\item Одна из самых популярных тестовых сцен в 3D графике
\item Модель атриума дворца Спонца в Дубровнике (Хорватия)
\end{itemize}
\end{frame}

\begin{frame}[fragile]
\frametitle{Сцена: формат OBJ}
Wavefront OBJ -- один из общепринятых форматов сцен
\begin{itemize}
\item Текстовый, достаточно простой для чтения
\item Содержит координаты вершин, нормали и текстурные координаты
\item Содержит индексы вершин, образующих треугольники
\item Может содержать много объектов
\item Может ссылаться на MTL-файл (Material Template Library), содержащий описания материалов
\item MTL может в свою очередь ссылаться на текстуры (альбедо, нормали, и т.п.)
\end{itemize}
\end{frame}

\begin{frame}[fragile]
\frametitle{Сцена: формат OBJ}
Что нужно из описания материалов:
\begin{itemize}
\item Текстура альбедо \verb|map_Ka| -- коэффициент для диффузного освещения (без нормализации; specular компоненту на неё \textbf{не умножаем})
\item Тестура прозрачности \verb|map_d| -- делаем discard, если в текстуре значение меньше 0.5 (нужно использовать в основном шейдере и в шейдере для shadow map), блендинг делать \textbf{не нужно}
\item Коэффициент отражения (glossiness) \verb|Ks| -- множитель для мощности specular компоненты
\item Показатель отражения (power) \verb|Ns| -- показатель степени для specular компоненты
\item Specular компоненту используем с нормализацией
\end{itemize}
\end{frame}

\begin{frame}[fragile]
\frametitle{Сцена: формат OBJ}
\begin{itemize}
\item Можно написать загрузчик руками или дописать тот, что есть у нас в практиках
\item Можно использовать какую-нибудь библиотеку, например Assimp или TinyOBJ
\item \textbf{N.B.} Есть два соглашения о том, как идут текстурные координаты: снизу вверх, или сверху вниз. Если текстуры будут выглядеть странно, попробуйте отразить текстурные координаты по оси Y: \verb|y = 1 - y|
\end{itemize}
\end{frame}

\begin{frame}[fragile]
\frametitle{Советы}
\begin{itemize}
\item Можно иметь один общий набор VAO + VBO + EBO, и для каждого объекта хранить только диапазон индексов
\item Можно иметь по VAO + VBO + EBO на каждый объект
\item Скорее всего, у вас будет 1 draw call (\verb|glDrawElements| или т.п.) на один объект
\item Напоминание: для рендеринга VBO и EBO не нужны, только \verb|glBindVertexArray| и \verb|glDraw*|!
\end{itemize}
\end{frame}

\begin{frame}[fragile]
\frametitle{Советы}
\begin{itemize}
\item В сумме у вас будет 2 shadow map (одна обычная, для направленного источника; вторая -- cubemap, для точечного)
\item В сумме у вас будет 8 texture view для shadow map (одна для направленного источника; одна для семплинга cubemap текстуры в основном шейдере; 6 для генерации граней cubemap текстуры)
\item Размер shadow map лучше взять побольше, если не будет тормозить (4096x4096)
\item Размывать тени можно прямо в результирующем шейдере, при чтении из shadow map, а можно -- в два прохода, сепарабельным фильтром за \begin{math}O(N)\end{math} \textit{(см. separable gaussian blur)}, за это будет больше баллов
\item Радиус размытия -- как можно больше, пока не начинает тормозить :)
\end{itemize}
\end{frame}

\begin{frame}[fragile]
\frametitle{Советы}
\begin{itemize}
\item При загрузке полезно посчитать bounding box сцены, чтобы потом по нему вычислять матрицы для теней
\item Лучше вычислять проекцию тени для направленного источника света используя \textbf{пересечение} bounding box'а видимой области камеры (см. конец лекции 4) и bounding box'а самой сцены, чтобы ничего лишнего не попало в shadow map, за это будет больше баллов (\textbf{NB:} для этого нужно включить \textit{depth clamping})
\item При инициализации wgpu\_app (и в С++, и в Rust) \textbf{нужно} передавать \mintinline{cpp}|srgb = false| (что это значит, мы узнаем позже)
\end{itemize}
\end{frame}

\begin{frame}[fragile]
\frametitle{Баллы}
\footnotesize
\begin{itemize}
\item \textbf{1 балл}: движется камера
\item \textbf{1 балл}: геометрия сцены загружается и рисуется
\item \textbf{1 балл}: есть текстуры альбедо
\item \textbf{1 балл}: есть текстуры прозрачности
\item \textbf{1 балл}: есть ambient освещение + два движущихся источника света
\item \textbf{2 балла}: есть тени от направленного источника света
\begin{itemize}
\item \textbf{+1 балл}: PCF + blur
\item либо \textbf{+1 балл}: VSM + blur
\item либо +\textbf{2 балла}: VSM + separable blur
\end{itemize}
\item \textbf{1 балл}: проеция тени для направленного источника вычисляется на основе camera frustum
\item \textbf{3 балла}: есть тени от точечного источника света
\begin{itemize}
\item \textbf{+1 балл}: PCF + blur
\item либо \textbf{+1 балл}: VSM + blur
\item либо \textbf{+2 балла}: VSM + separable blur
\end{itemize}
\end{itemize}
Всего: \textbf{15 баллов}

Защита заданий на практике \textbf{через 2 недели}.
\end{frame}

\end{document}
